\SourcePage{303}
\subsection{Lefschetz 束}\label{sec:19b}
设 $\overline X$ 为 $(n+1)$ 维光滑射影簇，带有丰沛线丛 $\mathcal L$。取两个线性无关截面 $s_0,s_\infty\in H^0(\mathcal L)$，考虑有理函数 $\overline p=s_0/s_\infty$，其纤维为 $z\in\CP^1$ 对应的超曲面 $\{s_0(y)=zs_\infty(y)\}$。$\overline p$ 的不定点集为基点集 $B$，由 $s_0$ 与 $s_\infty$ 同时为零定义。若 $B$ 光滑，且 $\overline p$ 只有非退化临界点，每个纤维至多含一个，就称 $\overline p$ 为 \emph{Lefschetz 束}。还假设 $s_\infty^{-1}(0)$ 本身光滑，并考虑 $\overline p$ 在其补集上的限制，记为
\begin{equation}\label{eq:19.1}
p:X=\overline X\setminus s_\infty^{-1}(0)\longrightarrow\CC.
\end{equation}
在 $\mathcal L$ 上选择一个给出 Kähler 形式的度量。如\stref{例题}{19.3}(i)，$s_\infty$ 的范数的负对数定义一个有限型多重次调和函数 $\phi$。所得设置与第\sr{sec:15}{15 节}引入的正合 Lefschetz 纤维化概念密切相关，但并不完全相同。有两种处理这一差别的方式：一是说明\eref{19.1}的几何性质足够好，可顺利推广所需论证；另一种是修改给定情形，使它完全符合前面建立的形式体系。

第一种方法的要点，是说明\eref{19.1}上由 $\omega_\phi$ 给出的联络具有良好的平行移动映射，尽管纤维并不紧。为此，沿基点集作爆破，得到 $\overline p$ 的图像
\[
\widehat X=\{(z,y)\in\CP^1\times\overline X:s_0(y)=zs_\infty(y)\}.
\]
它带有投影 $\widehat p:\widehat X\to\CP^1$ 与 $\widehat X\to\overline X$。通过第二个映射拉回给定 Kähler 形式，记所得形式为 $\omega_{\widehat X}$。它沿例外除子 $\widehat B=\CP^1\times B$ 退化，因而不再是 Kähler 形式，但在光滑纤维上的限制仍是 Kähler 形式，因此定义一个辛联络。该联络实际上在 $\widehat B$ 上平凡，所以在补集 $\widehat X\setminus\widehat B$ 上，其平行移动在奇点之外有良好定义。另一方面，从 $\widehat p:\widehat X\setminus\widehat B\to\CP^1$ 中去掉无穷远处的纤维，就恢复了 $p:X\to\CC$ 及由 $\omega_\phi$ 给出的联络。

这一论证还给出更多信息，在分次背景下尤其有意义。具体地，假设对某个 $d\in\Z$，有
\begin{equation}\label{eq:19.2}
\mathcal K_{\overline X}\cong\mathcal L^{\otimes-d}.
\end{equation}
在此同构下，$s_\infty^{\otimes-d}$ 的原像是一个有理全纯 $(n+1)$ 形式，它沿 $s_\infty^{-1}(0)$ 可能有零点，若 $d<0$，或极点，若 $d>0$。因此，它在 $X$ 上的限制为全纯体积形式 $\eta_X$。有了这一形式，准确地说，有了相关的相对二次形式 $\eta_X^2/dz^2$，全部平行移动映射都如\secref{15c}那样自然带有分次。考虑路径为充分大的圆周的情形。纤维化在无穷远处光滑延拓，但相对二次形式
\SourcePage{304}
不能如此延拓。因此，无穷远处的分次单值变换同痕于某个平移 $S^k$，其中 $k\in2\Z$；这里 $S^k$ 是带常值分次 $-k$ 的恒等映射。显然，$k$ 仅依赖于 $\eta_X/dz$ 沿无穷远纤维的极点阶数 $d-2$。为简便，改取平凡纤维化 $\CC^*\times Y\to\CC^*$，相对体积形式为 $z^{d-2}\eta_Y$。若 $\lambda(t)=e^{2\pi it}$ 为单位圆周，则沿 $\lambda|[0,t]$ 的平行移动是带分次 $(2d-4)t$ 的 $\Id_Y$。因此，沿圆周走完一圈后得到 $k=4-2d$。\znote{296}{极点阶 $d-2$ 的来历：(19.2) 给出 $\mathcal K_{\overline X}\cong\mathcal L^{\otimes-d}$，即 $\eta_X$ 沿除子 $s_\infty^{-1}(0)$ 有 $d$ 阶极点；而 $z=s_0/s_\infty$ 在该除子上恰有二阶极点，故相对形式 $\eta_X/dz$ 的极点阶为 $d-2$。取平凡化后相对形式写成 $z^{d-2}\eta_Y$，绕一圈时 $\arg z$ 增加 $2\pi$，分次增量即 $(d-2)\cdot2=2d-4$；又因 $S^k$ 携带的是常值分次 $-k$，故得 $k=4-2d$。这里的关键是：绕无穷远的单值变换只在分次意义下与恒等相差一个常值，因为决定分次的是相对二次形式而不是形式本身，故平移量必为 $2\Z$ 中的常数，且只能由沿无穷远纤维的极点阶读出。[来源：A，原书 (19.2) 及本段计算。]}

现在讨论第二种方法的一种具体实现，将\eref{19.1}变成第\sr{sec:15}{15 节}意义下的正合 Lefschetz 纤维化。需选择基点 $*\in\CC$；为简化记号，取 $*=0$，因此假设 $0$ 是 $p$ 的正则值。固定 $\CC$ 上的标准辛形式 $\omega_{\CC}$ 及其标准原始一形式 $\theta_{\CC}$。随后将对某个 $r\ge0$，把给定 $\omega_X$ 变形成 $\omega_X^r=\omega_X+rp^*\omega_{\CC}$，类似地把 $\theta_X$ 变形成 $\theta_X^r=\theta_X+rp^*\theta_{\CC}$。尽管缺少紧性，Moser 论证仍然适用，因为满足 $\omega^r(\cdot,Y^r)=\partial\theta_X^r/\partial r=p^*\theta_{\CC}$ 的向量场 $Y^r$，处处与底空间 $\CC$ 上向内的径向向量场 $-z\partial_z$ 的水平提升成正比例。所得结果是一族正合辛嵌入
\[
\rho^r:X\longrightarrow X,\qquad
(\rho^r)^*\omega_X^r=\omega_X,\qquad
(\rho^r)^*\theta_X^r=\theta_X+d(\text{某个函数}).
\]
记住这些嵌入的一个特殊性质很有用；它直接来自上述对 $Y^r$ 的描述：

\StatementNumber{6}\begin{lemma}\label{st:19.6}
取紧子集 $K\subset X$。则存在另一个紧子集 $\widetilde K\subset X$，使对所有 $r\ge0$，都有 $\rho^r(K)\subset\widetilde K$。
\end{lemma}

下一步是对给定辛联络，使用径向平行移动，为\eref{19.1}构造无穷远处的平凡化。它由开子集 $U\subset X$ 及 $U$ 与 $\CC\times U_0$ 之间的逐纤维微分同胚构成，其中 $p|(X\setminus U)$ 为固有映射，$U_0=U\cap X_0$。在此平凡化下，
\[
\theta_X=\theta_{X_0}+\kappa+dR,\qquad
\omega_X=\omega_{X_0}+d\kappa,
\]
其中 $\theta_{X_0},\omega_{X_0}$ 为 $\theta_X,\omega_X$ 在 $X_0$ 上的限制；$\kappa$ 是在纤维切向量上为零的一形式，$R$ 为函数。事实上，在所有 $z=0$ 的点 $(z,x)\in\CC\times U_0$ 上，二者都为零。现在希望修改平行移动，使它在给定平凡化下于无穷远处平凡。为此，取 $X_0$ 上紧支撑的函数 $l$，使它在 $U_0$ 之外等于 $1$，考虑修改后的形式
\begin{equation}\label{eq:19.3}
\widetilde\theta_X=\theta_{X_0}+l\kappa+d(lR)+rp^*\theta_{\CC},\qquad
\widetilde\omega_X=d\widetilde\theta_X,
\end{equation}
其中 $r\gg0$，并约定在平凡化适用的区域之外，将它们延拓为 $\theta_X^r,\omega_X^r$。取大圆盘 $S\subset\CC$，
\SourcePage{305}
赋予限制 $\theta_S=\theta_{\CC}|S$、$\omega_S=\omega_{\CC}|S$。“大”首先意味着 $S$ 的内部包含 $p$ 的全部临界值；稍后还会施加更多条件。取紧子集 $E_0\subset X_0$，使它为某个多重次调和函数的大次水平集，并考虑一块 $E\subset p^{-1}(S)$，使它在无穷远处形如 $S\times E_0$。这里约定，$E$ 包含 $p^{-1}(S)$ 中在无穷远平凡化下看不到的全部部分。此外，选择 $E_0$，使 $l$ 支撑于其内部。令 $\omega_E=\widetilde\omega_X|E$、$\theta_E=\widetilde\theta_X|E$，考虑限制 $\pi=p|E:E\to S$。最后，取与 $\omega_E$ 相容的殆复结构 $I_E$，使它在临界点附近与给定 $I_X$ 一致，使 $\pi$ 伪全纯，并在无穷远附近的平凡化下形如自然乘积结构 $i\times I_{X_0}$。我们断言，所得结果是正合 Lefschetz 纤维化。由于限制在纤维上时 $\widetilde\omega_X=\omega_X$，而且 $r$ 很大，显然 $\widetilde\omega_X$ 仍为辛形式，负 Liouville 向量场沿 $\partial^vE=\pi^{-1}(\partial S)$ 指向内部；考察\eref{15.6}时已见过这一论证。另一方面，已知负 Liouville 场沿 $E_0$ 的边界指向内部，因此由构造，它沿全部 $\partial^hE$ 也如此。事实上，已将纤维化安排成在 $\partial^hE$ 附近按\stref{笔记}{15.2}的意义平凡。$\partial^hE$ 的弱 $I_E$-凸性同样来自 $E_0$ 的相应性质；$\partial^vE$ 的对应性质则由将伪全纯曲线投影到底空间得到。最后，假设 $\overline X$ 带有同构\eref{19.2}。由构造，$I_E$ 作为殆复结构同伦于 $I_X|E$；沿此同伦传递 $\eta_X|E$，即可得到 $(E,I_E)$ 的复体积形式 $\eta_E$。

验证了所需性质后，还须考察\eref{19.3}对辛结构的改变及其对 Fukaya 范畴的影响。$X$ 上如\stref{例题}{19.3}(i)给出的自然多重次调和函数 $\phi$ 为有限型。取次水平集 $X^{\phi,c}$，其中 $c$ 大于全部临界值。将它作为\stref{引理}{19.6}中的 $K$，考虑相应的 $\widetilde K$。构造时，可使 $\widetilde\omega_X=\omega_X^r$ 与 $\widetilde\theta_X=\theta_X^r$ 在 $\widetilde K$ 上成立，于是存在 $X^{\phi,c}$ 到 $E$ 的正合辛嵌入。由\stref{命题}{10.5}，$\F(E)$ 因此含有一个拟等价于 $\F(X^{\phi,c})$ 的全 $\Ai$ 子范畴。这已完全满足我们的目的，即利用 Lefschetz 纤维化得出关于 $X$ 的结论；$\F(E)$ 中任何其他对象都可视为额外对象而忽略。若考察纤维化的几何，而不是总空间 $E$，事情会更简单。按定义，纤维 $E_0$ 是 $X_0$ 上某个多重次调和函数的大次水平集。$\pi:E\to S$ 的平行移动映射当然不是 $p:X\to\CC$ 的平行移动的限制，但在 $E\setminus\partial^hE$ 的一个大紧子集上，两者一致；这是无穷远平凡化没有改动的部分。可以将各项选择安排成使 $E_0$ 中每个 Lagrangian 子流形都能由负 Liouville 流推入该部分。特别地，这使我们能沿用前面关于整体单值变换的结果：若 $\lambda$ 为 $S\setminus\operatorname{Critv}(\pi)$ 中沿正向平行于边界的环路，$L$ 为纤维
\SourcePage{306}
$E_{\lambda(0)}$ 中的正合 Lagrangian 子流形，则单值变换像 $h_\lambda(L)$ 同痕于 $L$。存在\eref{19.2}时，还有更强的版本：任意正合 Lagrangian 膜 $L^\#$ 在 $h_\lambda$ 的自然分次版本下的像满足
\begin{equation}\label{eq:19.4}
h_\lambda(L^\#)\simeq S^{4-2d}L^\#.
\end{equation}

\StatementNumber{7}\begin{proposition}\label{st:19.7}
设 $\pi:E^{2n+2}\to S$ 为来自 Lefschetz 束的正合 Lefschetz 纤维化，带有由\eref{19.2}得到的相对复体积形式。固定基点 $*\in S$，令 $V^\#=(V_1^\#,\ldots,V_m^\#)$ 为纤维 $M=E_*$ 中的消失圈基，任意选择分次使它们成为 Lagrangian 膜；当 $n=1$ 且 $\operatorname{char}(\K)\ne2$ 时，取非平凡 Pin 结构。设 $\A=\Pi(\Tw\F(M))$ 为扭曲复形 $\Ai$ 范畴的分裂闭包，故按定义 $H^0(\A)=D^\pi\F(M)$。若 $d\ne2$，则 $V^\#$ 中的对象分裂生成 $\A$。
\end{proposition}
\begin{proof}
结合\eref{19.4}、\stref{推论}{17.17}及\stref{推论}{5.8}即可。\znote{298}{$d\ne2$ 正是 (19.4) 中平移量 $4-2d$ 非零的条件：$d=2$ 时平移为零，绕边界一圈的分次单值变换就是恒等映射，\eref{19.4} 不再提供任何信息；而“消失圈分裂生成”恰需这个非零平移来配合\stref{推论}{17.17}（非平凡平移给出沿环路单值变换的非平凡性，从而使相应对象在分裂闭包中生成）。下面的\stref{笔记}{19.8}用反例说明此时结论确实不成立，故该限制不是技术性的。[来源：A，原书 (19.4)、命题 19.7 的证明及\stref{笔记}{19.8}。]}
\end{proof}

\StatementNumber{8}\begin{note}\label{st:19.8}
为说明条件 $d\ne2$ 确实必需，取 $\overline X=\CP^1\times\CP^1$，$\mathcal L=\mathcal O_{\overline X}(1,1)$。此时 $X_0=\CC^*$，两个消失圈都同痕于同一圆周 $L=S^1\subset\CC^*$，带有非平凡 Pin 结构。将它记为 $L^\#$，将带平凡 Pin 结构的对应物记为 $L^\#\otimes\beta$。若系数域 $\K$ 的特征不为 $2$，则由\eref{12.14}，$HF^*(L^\#,L^\#\otimes\beta)=H^*(L;|\beta|_{\K})=0$。因此，$L^\#\otimes\beta$ 不在由 $L^\#$ 分裂生成的子范畴中。\znote{298}{这个例子可逐项核对：$\mathcal L=\mathcal O(1,1)$，而 $K_{\CP^1\times\CP^1}=\mathcal O(-2,-2)=\mathcal L^{\otimes-2}$，故 $d=2$。$D=s_\infty^{-1}(0)$ 是 $(1,1)$ 除子，$X=\overline X\setminus D$；纤维 $X_0$ 为 $s_0^{-1}(0)\setminus D$，其中 $s_0^{-1}(0)\cong\CP^1$ 与 $D$ 恰交于两点，故 $X_0=\CC^*$，两个消失圈都同痕于 $S^1\subset\CC^*$。$d=2$ 使 (19.4) 的平移为零，这正是结论失效的根源。再看后半句：$HF^*(L^\#,L^\#\otimes\beta)=H^*(L;|\beta|_{\K})$ 之所以为零，是因为 $\beta$ 使系数成为秩一的非平凡局部系统（沿 $S^1$ 的单值变换为 $-1$），此时 $H^0$ 与 $H^1$ 同时为零，故 $L^\#\otimes\beta$ 无法由 $L^\#$ 分裂生成。[来源：A，原书 (19.2)、\stref{例题}{19.3}(i)、\eref{12.14} 及本节。]}

尽管如此，仍可尝试将 $d=2$ 的一般理论再推进一些。写 $L_k=\tau_{V_{k+1}}\cdots\tau_{V_m}(L)$，赋予诱导的膜结构。在正合三角
\SourceDiagram{diagram-19-b-triangle}{0.384230355}
中，将所有 $\to$ 映射复合，得到数
\[
c_{L^\#}\in HF^0(L_m^\#,L_0^\#)\cong HF^0(L^\#,L^\#)\cong H^0(L;\K)=\K.
\]
考察\stref{推论}{5.8}的证明，可知 $c_{L^\#}=0$ 是 $L^\#$ 属于由消失圈分裂生成的子范畴的一个充分条件。另一方面，回顾\stref{推论}{17.18}中这些正合三角的构造，可知每个 $\to$ 映射都是计数给定 Lefschetz 纤维化一小块中的伪全纯截面的不变量，该小块仅含一个临界点。利用标准粘合定理，可将 $c_{L^\#}$ 写成整个纤维化的一个不变量 $\Phi_{E/S,*}$，其中移动的 Lagrangian 边界条件 $F$ 实现同痕
\SourcePage{307}
\eref{19.4}。等价地，从原来的代数几何数据看，可将 $F$ 视为 $X$ 中在 $p:X\to\CC$ 的底空间内一个大圆周上纤维化的 Lagrangian 子流形，因此它靠近无穷远纤维。于是，$c_{L^\#}$ 计数 $X$ 中边界位于 $F$ 的全纯圆盘。可尝试通过向法锥退化，或辛切割，进一步分析这一情形。该过程使 $X$ 退化为两块之并：$\overline X$ 与 $\mathcal L^{-1}\to s_\infty^{-1}(0)$ 的总空间，两者沿 $s_\infty^{-1}(0)$ 粘合。Lagrangian 子流形 $F$ 移入第二块，可用类似\bcite{79}的技术计数全纯圆盘。这里不再尝试作更具体的说明。
\end{note}
